Our goal is to design a new dataplane operating system architecture that enables scalable, microsecond-scale preemptive scheduling on modern multi-core servers.
Achieving this requires reconciling a fundamental tension: delivering the near bare-metal performance required for fine-grained scheduling, while ensuring rigorous system safety when exposing low-level hardware primitives (e.g., LAPIC timers) to untrusted userspace components.
We identify two major barriers to bridging this gap:

\paragraph{C1: Virtualization-based isolation incurs prohibitive overhead for microsecond-scale scheduling.}
Existing virtualization techniques (e.g., KVM with APICv) provide robust isolation between physical hardware and untrusted code.
However, the architectural cost of maintaining this isolation boundary is incompatible with microsecond-scale scheduling intervals.
Table \ref{tab:operation-overhead} quantifies the overhead of key preemption primitives on bare metal versus virtualization.
While hardware assists like APICv optimize interrupt delivery (i.e., avoiding \texttt{VMEXIT} and \texttt{VMENTER}), they fail to optimize privileged configuration operations.
For regular context switches, the overhead increases from 151 cycles on bare metal to over 6,400 cycles with APICv.
For timer arming, which happens at every context switch, the overhead is 490 cycles on bare metal but skyrockets to over 5,600 cycles even with APICv, due to the unavoidable \texttt{VMEXIT}/\texttt{VMENTER}.
For a system targeting a 5 $\mu$s timeslice (approx. 12,500 cycles), spending nearly 50\% of the quantum just to program the timer is unacceptable.
For interrupt delivery, APICv reduces the overhead to 2,717 cycles by posted interrupts, reaching near bare-metal performance.

\yihan{Need to update the table}
\begin{table}[h]
\centering
\begin{tabular}{lccc}
\toprule
\multirow{2}{*}{\textbf{Operation}} & \multirow{2}{*}{\textbf{Bare Metal}}  & \multicolumn{2}{c}{\textbf{Virtualization}} \\ \cmidrule(l){3-4} 
                                    & & \small with APICv     & \small w/o APICv     \\
% \midrule
Mode Switch                & 151            & 6,439          & 10,108            \\
Timer Setup and Fire (Immediate)                        & 3,496            & 5,539          & 10,854             \\
Timer Setup and Fire (EOI)                        & 3,496            & 435          & 4,381             \\
% Interrupt Delivery                 & 2,960          & 2,717          & 6,305             \\
\bottomrule
\end{tabular}
\caption{Latency (in cycles) of key preemption primitives on a 2.5GHz Intel Xeon Gold 6548Y (Linux 6.8).\textit{Context Switch} is proxied by a syscall (\texttt{getpid}) vs. a VMCALL.\textit{Timer Arming} compares \texttt{hrtimer\_start} (bare metal) vs. \texttt{wrmsr} to the virtual APIC.\textit{Interrupt Delivery} measures the latency from timer expiration to handler entry.}
\label{tab:operation-overhead}
\end{table}

\paragraph{C2: Bypassing virtualization is unsafe due to the lack hardware multiplexing mechanism for system-critical resources.}
While modern processors provide IOMMUs to safely expose PCIe devices (e.g., NICs, GPUs) to userspace, no equivalent hardware mechanism exists for core-local resources like the LAPIC.
The LAPIC is a privileged component integral to system stability; it manages not just user preemption, but also kernel watchdog timers and IPIs.
Consequently, standard OS kernels strictly prohibit userspace access to the LAPIC.
Bypassing this restriction to grant an untrusted userspace scheduler direct access creates severe safety risks: a malicious or buggy application could misconfigure the timer to starve the kernel, trigger interrupt storms, or soft-lock the CPU, compromising the integrity of the entire server.

\subsection{\kernel Approach}
To overcome these challenges, \kernel introduces a new isolation model that achieves the performance of bare metal with the safety of virtualization.
Our approach relies on two key architectural abstractions:

To address C1, \kernel introduces the \emph{Shadow Ring}, a lightweight software enclave within the application process that securely isolates the trusted scheduler from untrusted application code (the tenant).
Instead of relying on heavy-weight virtualization, \kernel constructs a secure enclave \textit{within} the process address space using Intel MPK.
We partition the application process into two hardware-enforced domains:
\begin{itemize}
\item \textbf{The Tenant Ring (Untrusted):} Contains the application logic, data, and the standard userspace library OS components.
\item \textbf{The Shadow Ring (Trusted):} Contains the \kernel scheduler (\schedcore), its private internal state, and the sensitive code paths required to interact with hardware.
\end{itemize}
The Shadow Ring runs at the same privilege level (Ring 3) as the application but is logically isolated by MPK permissions.
This eliminates the costly \texttt{VMEXIT} or ring transitions for scheduler invocation, allowing \kernel to perform context switches and timer updates at near-native speeds.

To address C2, \kernel implements a \emph{Double-Gate} security model that governs all interactions between the untrusted tenant, the trusted scheduler, and the kernel.
We focus on three possible misbehaviors that could compromise system safety when exposing the LAPIC timer to userspace:
\begin{enumerate*}[leftmargin=*,label=(\arabic*)]
\item \textit{Illegal Entry into the Shadow Ring:} The untrusted tenant could exploit memory corruption vulnerabilities (e.g., ROP attacks) to execute the \texttt{WRPKRU} instruction and gain unauthorized access to the Shadow Ring.
\item \textit{Unauthorized LAPIC Access:} The untrusted tenant could attempt to impose arbitrary configurations on the LAPIC timer by trapping into the kernel and forging timer requests.
\item \textit{Preemption Bypass:} The untrusted tenant could intercept or ignore timer interrupts, preventing the scheduler from regaining control.
\end{enumerate*}
To mitigate these risks, \kernel employs a two-tiered gatekeeping mechanism:
\begin{enumerate*}[leftmargin=*,label=(\roman*)]
\item \emph{The Meso-Gate (Lateral Protection):} Enforces strict control-flow integrity to prevent the untrusted tenant from illegally entering the Shadow Ring.
\item \emph{The Hyper-Gate (Vertical Protection):} Validates all LAPIC timer configuration requests to ensure they originate from the legitimate scheduler in the Shadow Ring.
\end{enumerate*}
Additionally, \kernel implements a \emph{Secure Upcall} mechanism to ensure that timer interrupts atomically redirect execution to the trusted scheduler, preventing preemption bypass.

Since this work is still in progress, we defer the detailed design and implementation of these mechanisms to future versions of this chapter.
We introduce the first prototype of \sysname and present preliminary evaluation results in the following sections.
